%\begin{tabularx}{0.8\linewidth}{|X|}\hline
%\textbf{Pour chaque question, recopier sur la copie son numéro et la réponse
%correspondante.}\\
%\textbf{Pour cette partie, aucune justification n’est demandée.}\\
%\textbf{Pour les questions a choix multiple, une seule réponse est exacte.}\\ \hline
%\end{tabularx}
%\end{center}

\medskip

\textbf{Question 1}

$A=\dfrac{8}{12}+\dfrac{9}{12}=\dfrac{17}{12}$
\medskip

\textbf{Question 2}

$45-45\times 0,10=45-4,5=40,5$ \euro{}
\medskip

\textbf{Question 3}

Les diagonales se coupent en leur milieux et sont de même longueur, donc \\
\ Réponse B : c'est un rectangle.

\medskip

\textbf{Question 4}

\begin{align*}
5x-15&=20\\
5x&=20+15\\
x&=\dfrac{35}{5}\\
x&=7
\end{align*}
\medskip

\textbf{Question 5}

L'abscisse du point $A$ est -4.\\
Les coordonnées du point $B$ sont (-2;-1).
\medskip

\textbf{Question 6}
1 ; 3 ; 3 ; 8 ; \fbox{11} ; 12 ; 12 ; 19 ; 25 

La médiane est 11.
\medskip

\textbf{Question 7}

$5\times\cos(60)$

\textbf{Question 8}

$3+8+7=18$ donc 387 est divisible par 3 et par 9. 
\medskip
